% =====================================================================
%  Propagation Mechanics of an Embedded Agent:
%  How Testimony, Divergence, and Inquiry Arise Without Being Represented
%
%  Self-contained. States its own axioms; re-derives every primitive it
%  uses; cites only standard, externally published mathematics. Every
%  object is a finite weighted graph or a construction on one.
% =====================================================================
\documentclass[11pt,twocolumn,a4paper]{article}

\usepackage[utf8]{inputenc}
\usepackage[T1]{fontenc}
\usepackage{lmodern}
\usepackage[margin=1in]{geometry}
\usepackage{amsmath,amssymb,amsthm,mathtools}
\usepackage{thmtools}
\usepackage{enumitem}
\usepackage{microtype}
\usepackage{booktabs}
\usepackage{graphicx}
\usepackage{caption}
\usepackage{hyperref}
\usepackage{cleveref}
\usepackage{xcolor}
\usepackage{tikz}
\usetikzlibrary{arrows.meta,positioning,calc}

\hypersetup{
  colorlinks=true, linkcolor=blue!50!black,
  citecolor=green!50!black, urlcolor=blue!60!black
}

% ---- theorem environments ----
\declaretheoremstyle[
  spaceabove=6pt, spacebelow=6pt,
  headfont=\bfseries, notefont=\mdseries, notebraces={(}{)},
  bodyfont=\itshape, postheadspace=1em
]{thmstyle}
\declaretheoremstyle[
  spaceabove=6pt, spacebelow=6pt,
  headfont=\bfseries, notefont=\mdseries, notebraces={(}{)},
  bodyfont=\normalfont, postheadspace=1em
]{defstyle}

\declaretheorem[style=thmstyle, name=Theorem, numberwithin=section]{theorem}
\declaretheorem[style=thmstyle, name=Proposition, sibling=theorem]{proposition}
\declaretheorem[style=thmstyle, name=Lemma, sibling=theorem]{lemma}
\declaretheorem[style=thmstyle, name=Corollary, sibling=theorem]{corollary}
\declaretheorem[style=defstyle, name=Definition, sibling=theorem]{definition}
\declaretheorem[style=defstyle, name=Axiom, sibling=theorem]{axiom}
\declaretheorem[style=defstyle, name=Remark, sibling=theorem]{remark}
\declaretheorem[style=defstyle, name=Example, sibling=theorem]{example}
\declaretheorem[style=defstyle, name=Principle, sibling=theorem]{principle}

\crefname{axiom}{Axiom}{Axioms}
\Crefname{axiom}{Axiom}{Axioms}

% ---- notation ----
\newcommand{\R}{\mathbb{R}}
\newcommand{\N}{\mathbb{N}}
\newcommand{\Agent}{\mathsf{A}}
\newcommand{\Gph}{\Gamma}
\newcommand{\V}{V}
\newcommand{\E}{E}
\newcommand{\wt}{w}
\newcommand{\med}{\mathsf{m}}
\newcommand{\flo}{\beta}
\newcommand{\cut}{\partial}
\newcommand{\co}{\complement}
\newcommand{\str}{\sigma}
\newcommand{\Walk}{W}
\newcommand{\rec}{M}
\newcommand{\acct}{\mathfrak{a}}   % an account
\newcommand{\Purp}{x^{\star}}
\newcommand{\Beh}{\mathcal{B}}       % disposition space
\newcommand{\drive}{\Phi}
\newcommand{\cost}{c}
\newcommand{\att}{\alpha}
\newcommand{\dem}{D}
\newcommand{\abs}[1]{\lvert#1\rvert}
\newcommand{\norm}[1]{\lVert#1\rVert}
\newcommand{\dist}{d}

\title{\textbf{Propagation Mechanics of an Embedded Agent:\\
How Testimony, Divergence, and Inquiry Arise\\
Without Being Represented}}

\author{Kundai Farai Sachikonye\\[2pt]
\small AIMe Registry for Artificial Intelligence\\
\small \texttt{kundai.sachikonye@bitspark.com}}
\date{\today}

\begin{document}
\maketitle

\begin{abstract}
\noindent
We study what an agent embedded in a finite world \emph{does} when it is
engaged, and we show that a family of behaviours an observer would call
\emph{testimony}, \emph{divergent report} (including what is ordinarily
named a lie), \emph{rumour}, and \emph{questioning} arise from the agent's
ordinary activity \emph{without the agent representing any of them}. We take
an agent to have a purpose it pursues and a positive cost on every act,
and to access its world only by \emph{search}: to recall or report anything
is to run a bounded walk over a finite graph of what it has individuated.
From four axioms---a positive floor on the cost of any distinction,
purpose-directed search, a positive floor on the cost of any act, and the
non-completability of the world---we prove, as graph and order theorems,
that: \textbf{(T1)} recalling and searching are one operation, so every
report an agent gives is freshly searched, never fetched; \textbf{(T2)} two
searches from the same start to the same conclusion traverse different
interiors and neither can reconstruct the other's route (\emph{path
opacity}), so an agent cannot certify its own account against any standard;
\textbf{(T3)} the same event registers a different result in agents with
different internal graphs, and each registration is a correct individuation
in its own graph, so divergence of report between agents is not error in
either; \textbf{(T4)} no agent can read the value another's report would
have to match, so the property an observer names ``the account is false''
is undefined \emph{inside} any agent and is only ever a relation an outside
party computes between two accounts---the acting agent holds no
representation of it; \textbf{(T5)} an account passed from agent to agent
acquires each relay's search as a composed factor, so its content drifts
multiplicatively along a chain with no relay altering it on purpose; and
\textbf{(T6)} because no report resolves to a held value, an exchange cannot
close, so a further act---which an observer names a question---is forced by
the residual gap and is selected by that gap, not by any represented intent
to inquire. We then locate the vocabulary: we prove that ``lie,''
``rumour,'' and ``question'' are predicates on \emph{relations between
accounts}, computable by an observer or a second agent, and that none of
them names any state or act of the agent that produced the account. The
agent contains no concept of deceiving, gossiping, or asking; the concepts
live entirely on the side of whoever watches. We give constructions, an
observer-side audit, and a numerical study on explicit agents. The reading
of the formal objects as testimony and inquiry is marked as such; no
theorem depends on it.
\end{abstract}

\vspace{0.5em}
\noindent\textbf{Keywords:} embedded agent; search; recognition--search
identity; path opacity; receiver-relative report; unextractable value;
catalytic drift; forced response; testimony; emergence without
representation; non-player character.

% =====================================================================
\section{Introduction}
\label{sec:intro}
% =====================================================================

\subsection{The question}

An agent placed in a world and engaged by another party must produce
something: a recollection, a report, an answer. We ask what it is
\emph{doing} when it does this, and we arrive at an observation that
reorganises the usual picture. Behaviours that an onlooker readily names---
the agent \emph{testified}, it \emph{lied}, it \emph{spread a rumour}, it
\emph{asked a question}---turn out not to be things the agent does
deliberately, nor things it needs any concept of in order to do. They are
names an onlooker attaches, after the fact, to the output of the agent's
ordinary activity. The agent that ``lied'' did nothing structurally
different from the agent that ``told the truth.'' The vocabulary is on the
onlooker's side. Testimony, lying, and communicative meaning have a large
philosophical literature~\cite{coady1992,mahon2016,grice1975,%
sperberwilson1986,davidson1973}; our contribution is to show these
predicates need no representation \emph{inside} the agent.

This paper makes that observation precise and proves it. We show that an
embedded agent, equipped only with a purpose and a cost on its acts, and
accessing its world only by bounded search, \emph{emits} testimony,
divergent report, rumour, and questioning \emph{without representing any of
them}---and that the very properties by which an onlooker classifies these
outputs are undefined inside the agent and are only ever relations computed
from outside.

\subsection{The claim: emergence without representation}

Our central claim is one of \emph{emergence}, deliberately not one of
necessity. We do not prove that such an agent must, in any given world,
produce a lie or a question; that would require assumptions about what the
world supplies. We prove the cleaner and more defensible thing: \emph{these
outcomes occur without the agent representing them}. Nowhere in the agent's
machinery is there a concept of deceiving, of gossiping, or of asking; yet
its outputs are correctly so classified by an observer. We then show
\emph{where} the classifying vocabulary can attach at all, and find---as a
result, not a premise---that it attaches only to relations between
accounts, never to any state or act of the agent that produced one.

\subsection{Method}

Every object below is a finite weighted graph or a construction on one, and
every result is a theorem from four stated axioms and standard
mathematics~\cite{bondymurty2008,diestel2017}.
We define an agent's world as a graph of what it has individuated, its
report as a walk on that graph, and its purpose and act-cost as in a
companion development which we restate here so the paper stands alone. The
deductions are combinatorial (cuts and walks), order-theoretic (monotone
histories), and use only elementary facts about minimum cuts, which are
computable exactly by maximum flow~\cite{fordfulkerson1956,menger1927,%
edmondskarp1972,cormen2009}. Where a tempting claim could not be
proved at this standard it has been removed; every interpretive reading is
confined to a remark on which no theorem depends.

\subsection{Organisation}

\Cref{sec:world} fixes the agent, its world-graph, its purpose and
act-cost, and states the axioms. \Cref{sec:searchknowing} proves that
recalling is searching (T1). \Cref{sec:opacity} proves path opacity and
that an agent cannot certify its own account (T2). \Cref{sec:divergence}
proves receiver-relative report and that divergence is not error (T3).
\Cref{sec:label} proves that falsity is undefined inside the agent and
locates it as an observer-side relation (T4). \Cref{sec:drift} proves
catalytic drift of a relayed account (T5). \Cref{sec:inquiry} proves that a
further act is forced by the residual gap (T6). \Cref{sec:vocabulary}
gathers the result: where each name lives. \Cref{sec:validation} reports a
numerical study. \Cref{sec:discussion} discusses scope.

% =====================================================================
\section{The agent and its world}
\label{sec:world}
% =====================================================================

We build the agent's world as a graph of what it has told apart, and its
report as a search over that graph.

\begin{definition}[World-graph; medium]
\label{def:world}
The \emph{world-graph} of an agent is a finite weighted graph
$\Gph=(\V,\E,\wt)$ with $\wt:\E\to\R_{>0}$, together with a distinguished
vertex $\med\in\V$, the \emph{medium}, adjacent to every other vertex. The
vertices $v\neq\med$ are the \emph{items} the agent has individuated (the
things it has told apart); an edge is a \emph{contact}, its weight the cost
of the separation distinguishing its endpoints. The graph is connected. The
medium is the single vertex standing for ``everything else,'' against which
each item is told apart---individuation against a background being the
classical problem of reference and identity~\cite{quine1960,putnam1975}.
\end{definition}

\begin{remark}[The world-graph is a finite window on an infinite whole]
\label{rem:infinitewhole}
The world-graph is not constructed by the agent or handed to it as input;
it is what the agent has \emph{individuated so far} of a whole that exceeds
it. The agent is a finite subgraph of an infinite whole, and the
world-graph $\Gph$ is the portion of that whole visible from the agent's
position---the items it has told apart, together with the contacts among
them. The medium $\med$ is not a separate object bolted on for bookkeeping:
it \emph{is} the rest of the infinite whole, always already present, the
undifferentiated ``everything else'' against which each item is told apart.
On this reading non-completability (\Cref{ax:noncompletable}) is not an
extra hypothesis about the world growing but a direct consequence of the
whole being genuinely infinite: no finite stage of individuation exhausts
it, so there is always a further item, and the medium never empties. This
grounds the world-graph in a foundational account of the agent as a proper
part of an infinite whole, and removes any appearance that $\Gph$ is an
arbitrary given.
\end{remark}

\begin{definition}[Cut; separation cost; resting cut]
\label{def:cut}
For $\varnothing\neq S\subsetneq\V$ the \emph{cut} is
$\cut(S)=\{\{u,v\}\in\E: u\in S,\ v\notin S\}$ with weight
$\wt(\cut(S))=\sum_{e\in\cut(S)}\wt(e)$. The \emph{separation cost} of an
item $v$ is the minimum weight of a cut placing $v$ with the medium
excluded,
\[
\str(v)=\min_{S\ni v,\ \med\notin S}\wt(\cut(S)),
\]
attained on finite $\V$; a minimising side is a \emph{resting cut} of $v$
and is the $v$--$\med$ minimum cut, computable by maximum flow via the
max-flow/min-cut theorem~\cite{fordfulkerson1956,menger1927}.
\end{definition}

We now restate the agent's purpose and act-cost, so the paper is
self-contained.

\begin{definition}[Disposition space; purpose]
\label{def:purpose}
The agent has a compact convex \emph{disposition space}
$\Beh\subseteq\R^n$ and a continuously differentiable, $\lambda$-strongly
convex \emph{drive} $\drive:\Beh\to\R$ ($\lambda>0$)~\cite{%
boydvandenberghe2004}. Its \emph{purpose} is the unique minimiser
$\Purp=\arg\min_{x\in\Beh}\drive(x)$, the attractor of the descent
$\dot x=-\nabla\drive(x)$~\cite{khalil2002}; the \emph{residual to purpose}
at $x$ is $\dem(x)=\drive(x)-\drive(\Purp)\ge0$.
\end{definition}

\begin{definition}[Act; act-cost]
\label{def:act}
An \emph{act} of the agent (a step of search, a committed report) draws from
a finite repertoire, each act carrying a \emph{cost} $\cost>0$. A bounded
agent commits only finitely many acts per unit time, in the sense of
bounded rationality and resource-limited agency~\cite{simon1955,simon1996,%
russellnorvig2020}.
\end{definition}

The four axioms follow. They are meant to hold of any bounded agent
embedded in a world it did not finish building.

\begin{axiom}[Distinction floor]
\label{ax:distfloor}
No separation is free: there is $\flo>0$ with $\wt(e)\ge\flo$ for every
contact. Hence $\str(v)\ge\flo>0$ for every item (a $v$--$\med$ cut is
nonempty by connectedness).
\end{axiom}

\begin{remark}[The distinction floor is not an independent axiom]
\label{rem:floor-not-independent}
We state \Cref{ax:distfloor} for convenience, but it is \emph{not}
independent of the others: it follows both from what a distinction
\emph{is} (\Cref{thm:floorforced} below, via boundary thickness) and,
quantitatively, from finiteness at each stage. We keep it as a named axiom
so later sections may cite the floor directly, but the reader should regard
it as a derived fact, demoted to the status of a lemma by
\Cref{thm:floorforced}. The genuinely primitive commitments are
purpose-directed search (\Cref{ax:purpsearch}), the act floor
(\Cref{ax:actfloor}), and non-completability (\Cref{ax:noncompletable}).
\end{remark}

\begin{axiom}[Purpose-directed search]
\label{ax:purpsearch}
The agent accesses its world only by search directed at its purpose: to
recall, report, or answer is to run a walk on $\Gph$ whose steps are
selected to reduce the residual $\dem$ toward $\Purp$
(\Cref{def:searchprocess} below). The agent has no direct read-out of any
item's value independent of such a walk.
\end{axiom}

\begin{axiom}[Act floor]
\label{ax:actfloor}
No act is free: there is a floor on act-cost, $\cost\ge\flo_{\mathrm{act}}
>0$, so the agent commits only finitely many acts per unit time.
\end{axiom}

\begin{axiom}[Non-completable world]
\label{ax:noncompletable}
The world is non-completable: for every finite stage of the agent's
world-graph there is a further item not yet individuated. This is an order
condition (no last stage), not a claim about cardinality.
\end{axiom}

\begin{theorem}[The floor is constitutive of distinction, and attained]
\label{thm:floorforced}
Every separation of the agent's world-graph carries strictly positive
weight, and the floor $\flo>0$ is attained:
\begin{enumerate}[label=\textup{(\roman*)},leftmargin=2.2em]
\item \textbf{Positive by boundary thickness.} A separation of zero cost is
a boundary of zero thickness, and a zero-thickness boundary is not a
boundary---it fails to separate. Hence every contact that actually tells its
endpoints apart has positive weight, and each item's resting cut, being a
genuine separation from the medium, has $\str(v)>0$.
\item \textbf{Attained by finiteness at each stage.} At any finite stage the
agent's world-graph has finitely many contacts, so
$\flo:=\min_{e}\wt(e)$ is a minimum over a finite set of positive reals,
hence attained and strictly positive; and $\str(v)\ge\flo>0$ for every item
since a $v$--$\med$ cut is nonempty by connectedness.
\end{enumerate}
Thus \Cref{ax:distfloor} is not an independent stipulation but a
consequence of what a distinction is (i) together with finiteness (ii).
\end{theorem}

\begin{proof}
(i) To hold two items apart at zero cost is to hold them apart by a boundary
across which nothing is expended to tell one side from the other; but a
boundary that costs nothing to cross does not divide its two sides---they are
not, at zero cost, told apart, hence not two. So any contact that genuinely
separates has positive weight. An item's resting cut is by
\Cref{def:cut} the minimum $v$--$\med$ cut, a genuine separation of $v$ from
``everything else,'' so it too has positive total weight: $\str(v)>0$.
(ii) The world-graph at any finite stage has a finite edge set (each edge an
individuated contact); the infimum of finitely many positive reals is a
minimum, attained and strictly positive, so $\flo:=\min_e\wt(e)>0$ is
well-defined at that stage and, being a minimum, is not a loose bound.
Connectedness makes every $v$--$\med$ cut nonempty, so
$\str(v)\ge\flo>0$. The two together give \Cref{ax:distfloor} as a derived
fact.
\end{proof}

\begin{remark}[What non-completability does, and does not, establish]
\label{rem:noncompletable-role}
It is worth separating two roles the reader might conflate. The
\emph{floor} $\flo>0$ is established by \Cref{thm:floorforced}---by what
distinction is, not by non-completability. What
\Cref{ax:noncompletable} does is different and equally load-bearing: it
guarantees the world-graph is always \emph{expandable}---there is always a
further item to individuate, so the agent's graph never closes into a finite
catalogue. This is why no individuation resolves to a residue-free point
(\Cref{sec:label}): not because each individuation ``costs something''
(that is the floor's job) but because the comparison against the whole never
terminates, there being always more world. The floor keeps each separation
positive; non-completability keeps the world open. Both are needed, and they
are different facts.
\end{remark}

\begin{remark}[Why the floor matters here]
\label{rem:floormatters}
The floor is the single fact the whole paper turns on. Because every
individuation costs at least $\flo$, no item is ever resolved to a
zero-content point (\Cref{sec:label}); because every act costs at least
$\flo_{\mathrm{act}}$, search is bounded and its route is a committed,
non-repeatable history (\Cref{sec:opacity}). Everything an observer later
names---testimony, divergence, drift, inquiry---is a consequence of the
gap the floor keeps open.
\end{remark}

% =====================================================================
\section{Recalling is searching: no report is fetched}
\label{sec:searchknowing}
% =====================================================================

We first show that, for an agent of this kind, there is no distinction
between recalling something and searching for it: to produce a report is
always to run a fresh search, never to fetch a stored value.

\begin{definition}[Query; search process; report]
\label{def:searchprocess}
A \emph{query} to the agent is a target item $x^\ast$ (what the engagement
is about). A \emph{search process} for $x^\ast$ from a seed item $v_0$ is a
$\dem$-descending walk $\Walk=(v_0,v_1,\dots,v_k)$ on $\Gph$ ending at a
resting cut that individuates the agent's answer. The \emph{report}
(equivalently the \emph{account}) $\acct(\Walk)$ is the terminal resting
cut together with the item it individuates: what the agent returns.
\end{definition}

\begin{figure*}[!htbp]
    \centering
    \includegraphics[width=\textwidth]{figures/panel_1.png}
    \caption{\textbf{The Floor and Recall-Is-Search: No Item Is Individuated
    for Free, and No Report Is Fetched.}
    (\textbf{A})~Separation cost $\sigma(v)$ on the linear $y$-axis versus
    world size in items on the linear $x$-axis ($3$ to $9$), every item of
    $80$ random worlds (steel-blue), with the floor line marked (crimson).
    Every separation cost, computed as an exact minimum cut, sits on or
    above the floor: the empirical content of \Cref{thm:floorforced}, that
    identifying an item carries an irreducible residual $\sigma(v)\ge\flo>0$.
    (\textbf{B})~Histogram of the normalised separation cost $\sigma(v)/\flo$
    on the linear $x$-axis with the unit line marked (crimson), counts on the
    linear $y$-axis (sea-green). The support begins at $1$ and the smallest
    observed value is $1.092$---the floor is close to tight, so the residual
    the whole paper turns on is real and small.
    (\textbf{C})~Report produced (none / yes) on the $y$-axis versus
    committed acts in the search on the linear $x$-axis ($0$ to $12$)
    (steel-blue step, positive region shaded sea-green), with the zero-act
    case ringed (crimson). A report appears only after at least one committed
    act, the empirical content of \Cref{thm:recsearch}: recalling is
    searching and a zero-act fetch returns nothing.
    (\textbf{D})~Three-dimensional realised-floor surface over world size in
    items ($x$-axis, $3$ to $9$) and graph density ($y$-axis, $0.2$ to
    $0.9$) (viridis colourmap), viewing angle elevation $24^{\circ}$ azimuth
    $-58^{\circ}$. The realised floor stays strictly positive across the
    grid: no configuration individuates any item for free.}
    \label{fig:mechanics-1}
\end{figure*}

\begin{theorem}[Recognition--search identity]
\label{thm:recsearch}
For an agent accessing its world only by search (\Cref{ax:purpsearch}),
recalling an item and searching for it are one operation: producing the
report for a query $x^\ast$ is running a search process for $x^\ast$, and
there is no distinct ``fetch'' that returns a stored value without a walk.
\end{theorem}

\begin{proof}
By \Cref{ax:purpsearch} the agent has no read-out of an item's value
independent of a $\dem$-descending walk; the only operation that yields a
report is such a walk (\Cref{def:searchprocess}). A ``fetch'' would be a
value returned with no walk, i.e. a zero-act access; by \Cref{ax:actfloor}
every act costs $\ge\flo_{\mathrm{act}}>0$, so a zero-act access returns
nothing. Hence every report is the terminus of a search, and recalling is
searching.
\end{proof}

That proof, resting on \Cref{ax:purpsearch}, shows only that search-access
is what this agent has. The substantive question is \emph{why} fetch-access
is not merely absent but incoherent for an agent of this kind---why it could
not be added. The next theorem answers it, and does not depend on
\Cref{ax:purpsearch}.

\begin{theorem}[Fetch-access is architecturally incoherent under a monotone
history]
\label{thm:fetchincoherent}
Suppose, contrary to \Cref{ax:purpsearch}, the agent could fetch: return a
queried item's value as it was stored at some earlier stage, with no walk on
the current graph. Then the returned value is incoherent with the agent that
returns it. For the world-graph changes with every committed act---each act
costs $\ge\flo_{\mathrm{act}}>0$ (\Cref{ax:actfloor}) and increments the
monotone committed record, which never decreases (\Cref{cor:noreconstruct}),
so at the moment of the fetch the graph $\Gph$ is strictly later than the
graph at which the value was stored. A stored value is an individuation against the
\emph{earlier} graph and its medium; returned now, it is an answer
appropriate to a \emph{prior} agent---one the current agent cannot re-occupy,
since the record does not decrease (monotone non-return,
\Cref{cor:noreconstruct}). Fetch therefore delivers an account the present
agent could not itself have produced and cannot own; it is architecturally
incoherent for an agent with a monotone history, not merely excluded by
fiat.
\end{theorem}

\begin{proof}
Each committed act increments the monotone record (\Cref{ax:actfloor},
\Cref{cor:noreconstruct}) and alters the state in which any subsequent
report is produced; so the graph and history at the fetch strictly succeed
those at storage. A fetched value $\acct_0$ was the terminus of an individuation
against the stored-stage graph $\Gph_0$ and its medium $\med_0$. To be the
current agent's own account it would have to be an individuation against the
current $\Gph$ and $\med$; but $\Gph\neq\Gph_0$ (the record has grown) and
by \Cref{cor:noreconstruct} the agent cannot return to the stage of
$\Gph_0$. Hence $\acct_0$ is the account of a prior, now-unreachable agent,
not of the one holding it: fetch returns a value the present agent can
neither have generated nor re-occupy the standpoint of. The incoherence is
structural---it follows from monotone non-return alone---so fetch cannot be
coherently added to an agent with a growing record, independently of
\Cref{ax:purpsearch}.
\end{proof}

\begin{corollary}[Every report is freshly produced]
\label{cor:fresh}
No two engagements return a pre-stored answer; each runs its own search and
produces its report anew. In particular the agent's committed history grows
by at least one act per report (\Cref{ax:actfloor}), so the state in which a
report is produced is never exactly a prior state.
\end{corollary}

\begin{proof}
By \Cref{thm:recsearch} each report is a search terminus, and by
\Cref{ax:actfloor} each search commits $\ge1$ act of positive cost,
incrementing the agent's committed act-count. Two engagements therefore
occur at distinct act-counts, so their producing states differ.
\end{proof}

\begin{remark}[Interpretation: no canned lines]
\label{rem:nocanned}
\Cref{thm:recsearch} is the formal content of the intuition that a
convincing embedded agent does not read from a script. What it says when
engaged is searched, not fetched; and because the search runs against the
agent's current world-graph and purpose, the same query at two moments need
not return the same account. We use only the graph statement.
\end{remark}

% =====================================================================
\section{Path opacity: an agent cannot certify its own account}
\label{sec:opacity}
% =====================================================================

We now show that two searches reaching the same conclusion take different
routes, and that neither the endpoints nor any agent can recover which route
was taken. The consequence is that an agent has no access to a
``how I got here'' that it could check against a standard---it cannot
certify its own account.

\begin{definition}[Committed record; interior of a search]
\label{def:record}
The \emph{committed record} of a search $\Walk=(v_0,\dots,v_k)$ is
$\rec(\Walk)=k$, the number of steps committed. The \emph{interior} of
$\Walk$ is the sequence of intermediate items $(v_1,\dots,v_{k-1})$ it
passed through, excluding the seed and terminus.
\end{definition}

\begin{definition}[Agent state: terminus and count, not the walk]
\label{def:agentstate}
The agent's \emph{state} is the pair $(v_k,\rec)$ of its current item (the
terminus of its last search) and its committed record $\rec$. The state
carries the \emph{terminus} and the \emph{count}; it does not carry the
interior $(v_1,\dots,v_{k-1})$ of the walk that produced the terminus. An
agent retains where it arrived and how many acts it has committed, not the
route it took to arrive.
\end{definition}

\begin{remark}[Why the state excludes the interior]
\label{rem:state-excludes}
This is a substantive premise, made explicit: were the agent to store its
full walk history, it could in principle recover its route, and the
no-self-certification corollary below would fail. The premise is warranted
by the act floor: retaining each of the $k$ interior items would itself be
$k$ further committed acts of positive cost, growing without bound as
searches accumulate; an agent with a finite act-budget per unit time
(\Cref{ax:actfloor}) cannot carry the interior of every past search. The
economical and forced state is the terminus-plus-count of
\Cref{def:agentstate}, and the results of this section rest on it.
\end{remark}

\begin{theorem}[Path opacity]
\label{thm:opacity}
Let $\Walk_1,\Walk_2$ be two search processes with the same seed $v_0$ and
the same terminal report $\acct$. Then:
\begin{enumerate}[label=\textup{(\roman*)},leftmargin=2.2em]
\item their interiors may differ, and in a connected world-graph with any
item of degree $\ge2$ generically do; and
\item the report $\acct$ (seed and terminus) does not determine the
interior: from $\acct$ alone one cannot recover which of $\Walk_1,\Walk_2$
was run.
\end{enumerate}
\end{theorem}

\begin{proof}
(i) If some item on a $v_0$--terminus route has degree $\ge2$, there are two
admissible continuations from it (two contacts), yielding two walks with the
same endpoints and different interiors; connectedness supplies at least one
such branching item whenever the graph is not a bare path. (ii) The report
$\acct$ records only the seed and the terminal resting cut; the interiors of
$\Walk_1,\Walk_2$ differ (i) while their $\acct$ coincide, so the map
interior $\mapsto\acct$ is not injective and has no inverse: $\acct$ does
not determine the interior.
\end{proof}

\begin{corollary}[An agent cannot reconstruct its own route]
\label{cor:noreconstruct}
The committed record grows monotonically (each step adds one and none is
removed, since undoing a step is itself a further step), so once a search
terminates, its interior is a past state the agent does not re-occupy. The
agent holds its report $\acct$, not the walk that produced it; by
\Cref{thm:opacity}(ii) $\acct$ does not encode the interior. Hence the agent
cannot reconstruct the route by which it reached its own account.
\end{corollary}

\begin{proof}
Monotonicity of the record: a step increments it; an ``undo'' is a further
step incrementing it again, so it never decreases and the interior is not
re-entered. By \Cref{def:agentstate} the agent's state is the
terminus-plus-count $(v_k,\rec)$ and does \emph{not} include the interior;
\Cref{thm:opacity}(ii) shows the retained report $\acct$ underdetermines the
interior. With the interior neither carried in the state (\Cref{def:agentstate},
\Cref{rem:state-excludes}) nor recoverable from $\acct$, the route is
unavailable to the agent.
\end{proof}

\begin{figure*}[!htbp]
    \centering
    \includegraphics[width=\textwidth]{figures/panel_2.png}
    \caption{\textbf{Path Opacity and the Monotone Record: An Agent Cannot
    Reconstruct the Route to Its Own Account.}
    (\textbf{A})~Distinct interiors between the same endpoints on the linear
    $y$-axis versus world density on the linear $x$-axis ($0.3$ to $0.95$)
    (steel-blue), with the single-interior line marked (crimson, dashed).
    The count of distinct routes to one terminus rises above one as density
    grows, the empirical content of \Cref{thm:opacity}(i): same-endpoint
    searches generically differ in interior.
    (\textbf{B})~Number of searches on the linear $y$-axis versus interior
    length for a fixed seed and terminus on the linear $x$-axis (sea-green).
    Many interior lengths reach the same endpoints, the empirical content of
    \Cref{thm:opacity}(ii): the report does not determine the interior.
    (\textbf{C})~Three-dimensional bundle of up to four same-endpoint search
    walks, items placed on a ring in $(x,y)$ with search step on the
    $z$-axis (steel-blue, sea-green, amber, violet) from the common seed
    (crimson), viewing angle elevation $22^{\circ}$ azimuth $-60^{\circ}$.
    The walks share endpoints but trace different interiors, the geometric
    picture of path opacity.
    (\textbf{D})~Position on the linear $y$-axis versus the strictly
    increasing committed record $M$ on the linear $x$-axis ($1$ to $45$),
    points coloured by record (plasma), with the revisits of the initial
    position ringed (crimson). A position recurs but the record never
    returns, the empirical content of \Cref{cor:noreconstruct}: the
    committed record is monotone and the interior is a past state the agent
    does not re-occupy.}
    \label{fig:mechanics-2}
\end{figure*}

\begin{corollary}[No self-certification]
\label{cor:nocertify}
An agent cannot certify its own account against any standard that would
require its route: it holds the report but not the search that produced it,
and cannot recover the search. Whatever an observer might mean by ``the
agent knows how it reached this,'' the agent does not.
\end{corollary}

\begin{proof}
Certifying a route requires possessing it; by \Cref{cor:noreconstruct} the
agent does not possess and cannot recover its route. Hence no
route-dependent certification is available to the agent.
\end{proof}

\begin{remark}[Why this is the crux of ``definition-free'']
\label{rem:crux}
\Cref{cor:nocertify} is the technical heart of the paper's thesis. A concept
of lying would require the agent to compare its account against how things
stand and note a divergence it is responsible for. But the agent cannot even
recover its own route (\Cref{cor:noreconstruct}), let alone measure it
against a standard it does not possess (\Cref{sec:label}). The machinery an
agent would need in order to \emph{represent} its account as false is simply
absent. This is not a design omission; it is forced by path opacity and the
monotone record.
\end{remark}

% =====================================================================
\section{Divergent report is not error in either agent}
\label{sec:divergence}
% =====================================================================

Two agents engaged about the same event return different accounts. We show
this is not a fault in either: each account is a correct individuation in
its producer's own world-graph, and there is no shared value of which the
two are competing approximations.

\begin{definition}[Registration]
\label{def:registration}
The \emph{registration} of an external event in an agent is the resting cut
the event lands in within that agent's world-graph: its individuation
against \emph{that} agent's medium (its own ``everything else'').
\end{definition}

\begin{theorem}[Receiver-relative report]
\label{thm:receiverrel}
The same external event registers different resting cuts in agents with
different world-graphs, and each registration is a correct individuation at
its own agent's floor. No registration is privileged as the account; there
is no shared point-value of which the registrations are approximations.
\end{theorem}

\begin{proof}
The registration is the event's resting cut in the agent's graph
(\Cref{def:registration}); different graphs give different cuts (a sparse
graph individuates the event against few contrasts, a dense graph against
many), so the registrations differ. Each is the minimum $v$--$\med$ cut in
its own graph, hence a correct floor-attaining individuation there. Were
there a privileged shared value, a registration differing from it would be
incorrect; but each registration is the correct individuation in its own
graph, and---since completing an individuation to a residue-free point would
require exhausting the non-completable world
(\Cref{ax:noncompletable}, \Cref{rem:noncompletable-role})---no such shared
point-value exists. Hence none is privileged.
\end{proof}

\begin{corollary}[Divergence is not error]
\label{cor:divergence}
Two agents' accounts of one event may differ while each is a correct
individuation in its own world-graph. The divergence is not a mistake in
either agent; it is the difference of their world-graphs, and neither
account is the account.
\end{corollary}

\begin{proof}
Immediate from \Cref{thm:receiverrel}: each registration is correct in its
own graph, and there is no privileged value relative to which either could
be wrong.
\end{proof}

\begin{remark}[Two distinct claims, and which is earned here]
\label{rem:honestdiv}
It is important to separate two claims that are easy to run together.
\emph{First}: divergence between two honest agents' accounts is not error in
either. This is \Cref{cor:divergence}, and it is fully earned by
\Cref{thm:receiverrel} (T3)---each account is correct in its own graph, with
no privileged value relative to which either is wrong. \emph{Second}: a true
account and a false account are \emph{indistinguishable} in the producing
agent's state. This is a strictly stronger claim, and it is \emph{not}
established by T3 alone: it requires that falsity itself is undefined inside
the agent, which is proved only in the next section
(\Cref{thm:falsityundefined}, T4). We flag the dependency rather than assert
the indistinguishability here: the first claim stands now; the second is
deferred to \Cref{sec:label}, where it is earned. The two should not be
conflated---T3 gives ``divergence is not error,'' and only T4 gives ``true
and false are indistinguishable inside the agent.''
\end{remark}

\begin{figure*}[!htbp]
    \centering
    \includegraphics[width=\textwidth]{figures/panel_3.png}
    \caption{\textbf{Receiver-Relative Divergence and Falsity as a Relation:
    The Same Event, Two Correct Accounts.}
    (\textbf{A})~Dense-agent separation cost $\sigma(v)$ on the linear
    $y$-axis versus sparse-agent $\sigma(v)$ on the linear $x$-axis for the
    same event, $150$ agent pairs (steel-blue), with the identity diagonal
    marked (crimson, dashed). The points sit off the diagonal: the same event
    registers a different---yet, as an exact min-cut, correct---resting cut
    in each agent, the empirical content of \Cref{thm:receiverrel}.
    (\textbf{B})~Histogram of the registration difference across agents on
    the linear $x$-axis with zero marked (crimson), counts on the linear
    $y$-axis (amber). The mass sits away from zero: divergence is broad and
    is not error in either agent (\Cref{cor:divergence}).
    (\textbf{C})~Agreement matrix of resting-cut sides, account agent
    ($y$-axis) against elected standard agent ($x$-axis), eight agents, cells
    equal (green) or unequal (pale). Whether two accounts ``agree'' depends
    on which is taken as standard, the empirical content of
    \Cref{thm:falsityrelation}: falsity is a relation on a pair, not a state
    of the producer.
    (\textbf{D})~Three-dimensional surface of an event's separation cost over
    two agents' densities (both density axes $0.2$ to $0.9$) (viridis
    colourmap), viewing angle elevation $24^{\circ}$ azimuth $-58^{\circ}$.
    The cost varies smoothly with the receiving graph: there is no privileged
    shared value of which the registrations are approximations.}
    \label{fig:mechanics-3}
\end{figure*}

% =====================================================================
\section{Where ``false'' lives: an observer-side relation, not an agent-state}
\label{sec:label}
% =====================================================================

We now show that the property by which an observer classifies an account as
false is undefined inside any agent, and we locate it---by letting the proof
decide---as a relation an outside party computes between two accounts.

\begin{definition}[Falsity-predicate, provisionally]
\label{def:falsity}
Say an account $\acct$ is \emph{false relative to a standard} $\acct^\ast$
if $\acct$ and $\acct^\ast$ individuate different resting cuts of the same
queried item---i.e. they disagree as individuations. The \emph{standard}
$\acct^\ast$ is another account the classifier elects to treat as correct:
some other agent's account, or a consensus of accounts. It is \emph{always
itself a receiver-relative registration}, never a privileged view of the
world.
\end{definition}

\begin{remark}[There is no privileged ``world-state read-out'']
\label{rem:no-worldreadout}
One might be tempted to let the standard $\acct^\ast$ be a ``world-state
read-out''---the way things stand, read off directly. But the framework
forbids this: every account is an individuation against \emph{some} agent's
medium (\Cref{thm:receiverrel}), subject to the same floor and the same
receiver-relativity as every other. A purported world-state read-out would
be a registration not subject to those constraints---a view from nowhere, of
which no agent's account is a mere approximation---and \Cref{thm:receiverrel}
shows no such privileged value exists. So the standard is always another
agent's account, carrying no privilege over the account it judges; this is
exactly what \Cref{thm:falsityrelation} establishes---falsity is a symmetric
relation on a pair, not a comparison against a distinguished truth. We
therefore exclude ``world-state read-out'' from the admissible standards,
keeping the clean conclusion that falsity is never absolute.
\end{remark}

\begin{theorem}[Falsity is undefined inside the agent]
\label{thm:falsityundefined}
For the agent that produced $\acct$, the falsity-predicate of
\Cref{def:falsity} has no value: the agent holds no standard $\acct^\ast$
against which to evaluate it. Specifically,
\begin{enumerate}[label=\textup{(\roman*)},leftmargin=2.2em]
\item the agent holds no residue-free value of the queried item to serve as
$\acct^\ast$ (\Cref{ax:noncompletable,rem:noncompletable-role,thm:receiverrel}:
none exists and none is registered);
\item the agent cannot supply its own route as a standard, being unable to
recover it (\Cref{cor:nocertify}); and
\item any other agent's account is, to this agent, just another
receiver-relative registration with no privilege (\Cref{thm:receiverrel}),
not a standard.
\end{enumerate}
Hence ``$\acct$ is false'' evaluates to no truth-value using only the
producing agent's state: the predicate is undefined inside the agent.
\end{theorem}

\begin{proof}
The predicate (\Cref{def:falsity}) requires a standard $\acct^\ast$. (i) A
shared point-value does not exist (\Cref{ax:noncompletable,rem:noncompletable-role})
and is not registered (\Cref{thm:receiverrel}), so the agent has none to use. (ii) The
agent cannot recover its own route (\Cref{cor:nocertify}), so cannot use it
as $\acct^\ast$. (iii) Another agent's account carries no privilege over the
producer's own (\Cref{thm:receiverrel}), so does not function as a standard
within the producer. With no $\acct^\ast$ available in the agent's state,
the predicate has no argument to evaluate against and is undefined there.
\end{proof}

\begin{theorem}[Falsity is a relation an observer computes]
\label{thm:falsityrelation}
The falsity-predicate acquires a value only for a party that holds two
accounts and a choice of which is the standard: given $\acct$ and
$\acct^\ast$, an observer (or a second agent electing $\acct^\ast$ as
standard) can compute whether they individuate different resting cuts, a
decidable comparison of two cuts. The predicate is therefore a relation on
the pair $(\acct,\acct^\ast)$, evaluated outside the producing agent, and
its value depends on the classifier's choice of standard.
\end{theorem}

\begin{proof}
Given both accounts, ``they individuate different resting cuts of the same
item'' is a comparison of two finite edge sets, decidable directly. The
comparison requires holding both $\acct$ and $\acct^\ast$ and designating
one as standard---data a classifier has and the producing agent does not
(\Cref{thm:falsityundefined}). Hence the predicate is well-defined exactly
as a relation on the pair, computed by the classifier, and varies with the
choice of $\acct^\ast$.
\end{proof}

\begin{corollary}[The agent that ``lied'' represented no lie]
\label{cor:noliar}
An agent whose account an observer classifies as false performed no act
distinguishable, in the agent's own state or history, from one whose account
the observer classifies as true: both are search termini
(\Cref{thm:recsearch}), route-opaque (\Cref{cor:nocertify}), and correct in
the agent's own graph (\Cref{thm:receiverrel}). The property ``lie'' is
attached by the observer to the relation between accounts
(\Cref{thm:falsityrelation}); it names nothing the agent did or represented.
\end{corollary}

\begin{proof}
By \Cref{thm:recsearch,cor:nocertify,thm:receiverrel} the producing act is a
route-opaque search terminus, correct in the producer's graph, whether or
not an observer later calls it false; nothing in the agent's state carries
the falsity-predicate (\Cref{thm:falsityundefined}). The predicate is the
observer-side relation of \Cref{thm:falsityrelation}. Hence ``lie'' names
the relation, not the agent's act.
\end{proof}

\begin{principle}[To lie, an agent need not know it lies---because there is
nothing there to know]
\label{prin:nolie}
The classical intuition that lying requires knowing one lies is not merely
relaxed here; it is dissolved. There is no agent-internal fact of lying to
be known: the falsity-predicate is undefined inside the agent
(\Cref{thm:falsityundefined}) and exists only as a relation an observer
computes (\Cref{thm:falsityrelation}). An agent produces an account an
observer calls a lie by the same mechanism it produces one the observer
calls true. The knowing is not absent from the agent; the \emph{object} of
the knowing is.
\end{principle}

\begin{remark}[Interpretation: deception without a deceiver]
\label{rem:deception}
An onlooker watching such agents will correctly and usefully speak of lies,
honest reports, and mistaken testimony---these are real distinctions on the
onlooker's side, computed from relations between accounts. What the results
show is that the corresponding \emph{acts} are indistinguishable on the
agent's side. A synthetic world populated by these agents therefore exhibits
deception without any agent that deceives, in the precise sense that the
deception-predicate is real as an observer-relation and vacuous as an
agent-state. We use only the theorems; the reading is orientation.
\end{remark}

% =====================================================================
\section{Relayed accounts drift multiplicatively}
\label{sec:drift}
% =====================================================================

We now let an account pass from agent to agent and show that it changes
along the way---each relay composing its own search as a factor---with no
relay altering it on purpose. What an observer names rumour is this
composition.

\begin{definition}[Relay; search gain]
\label{def:relay}
A \emph{relay} of an account by agent $i$ is: agent $i$ receives an account
as a seed, runs its own search process (\Cref{def:searchprocess}) on its own
world-graph, and emits the resulting account. The \emph{search gain}
$g_i\in[0,1]$ of the relay is the fraction of the seed's individuating
contrasts that survive agent $i$'s search unchanged (so $g_i=1$ is a
verbatim pass, $g_i<1$ a partial reshaping); reading the surviving
contrasts as retained information places the drift in the setting
of~\cite{shannon1948,coverthomas2006}.
\end{definition}

\begin{theorem}[Catalytic drift along a chain]
\label{thm:drift}
Let an account pass through relays $1,2,\dots,n$, each with search gain
$g_i$. Say the reshapings are \emph{independent} when each relay's search
acts on the incoming contrasts through a factor unaffected by which
contrasts earlier relays reshaped---which holds, in particular, when the
relays' world-graphs are pairwise disjoint, or overlap so little that
correlations between successive reshapings are negligible. Under independent
reshapings the fraction of the original account's individuating contrasts
surviving the whole chain is the product
\[
G_n=\prod_{i=1}^{n}g_i,
\]
non-increasing in $n$ and strictly decreasing whenever a relay has $g_i<1$.
An account relayed through such a chain drifts multiplicatively away from its
origin, with no relay changing it on purpose. When the relays' world-graphs
overlap substantially the product is an \emph{approximation}: the true
surviving fraction is $G_n$ up to a correction that grows with the degree of
overlap, and $G_n$ is exact in the disjoint limit.
\end{theorem}

\begin{proof}
Each relay applies its search to the account it receives, preserving a
fraction $g_i$ of the incoming individuating contrasts (\Cref{def:relay}).
When the reshapings are independent, relay $i{+}1$'s preserved fraction is
$g_{i+1}$ regardless of which contrasts relay $i$ dropped, so the fractions
compose by multiplication and the surviving fraction after $n$ relays is
$G_n=\prod_{i=1}^n g_i$. Since each $g_i\le1$, $G_n$ is non-increasing in
$n$, strictly decreasing at any relay with $g_i<1$. If instead relay
$i{+}1$'s graph overlaps relay $i$'s, its search is constrained by the same
world-structure that constrained relay $i$, so the reshapings are
correlated and the true surviving fraction departs from the bare product by
a correction determined by the overlap; as the overlap tends to zero the
correction vanishes and $G_n$ is recovered exactly. In every case no relay's
search is directed at altering the account (each merely runs its own
purpose-directed search, \Cref{ax:purpsearch}); the drift is the composition
of ordinary searches, exact when independent and approximate otherwise.
\end{proof}

\begin{corollary}[Rumour is composed drift, not authored distortion]
\label{cor:rumour}
The divergence between an account's origin and its form after $n$ relays is
$1-G_n$, produced entirely by the composition of the relays' searches. No
relay authored the distortion; each ran the same search it runs for any
query. What an observer names a rumour is the multiplicative drift of
\Cref{thm:drift}.
\end{corollary}

\begin{proof}
Immediate from \Cref{thm:drift}: the surviving fraction is $G_n$, the drift
$1-G_n$, and each factor arises from an ordinary relay search
(\Cref{ax:purpsearch,def:relay}).
\end{proof}

\begin{remark}[Why relayed accounts must move]
\label{rem:mustmove}
A verbatim relay ($g_i=1$ for all $i$) would require each agent to pass the
account through unchanged---but by \Cref{thm:recsearch} a relay is a search
on the relay's own world-graph, which is in general a different graph
(\Cref{thm:receiverrel}), so the account is re-individuated and generically
$g_i<1$. Faithful transmission is the non-generic case; drift is the norm.
This is the mechanism by which a claim mutates as it crosses a population,
with every carrier honest by the lights of \Cref{sec:label}.
\end{remark}

\begin{figure*}[!htbp]
    \centering
    \includegraphics[width=\textwidth]{figures/panel_4.png}
    \caption{\textbf{Catalytic Drift and Forced Inquiry: Rumour Without an
    Author, a Question Without Intent.}
    (\textbf{A})~Surviving fraction $\prod_i g_i$ on the linear $y$-axis
    versus number of relays $n$ on the linear $x$-axis ($1$ to $14$), for
    three per-relay gains $g_i\in\{0.6,0.75,0.9\}$ (steel-blue, sea-green,
    amber). Each curve decays geometrically, the empirical content of
    \Cref{thm:drift}: a relayed account drifts multiplicatively along a
    chain.
    (\textbf{B})~Histogram of account drift $1-\prod_i g_i$ on the linear
    $x$-axis over $400$ random relay chains, counts on the linear $y$-axis
    (violet). Drift accumulates across the population, the empirical content
    of \Cref{cor:rumour}: rumour is composed drift, authored by no relay.
    (\textbf{C})~Residual gap $D$ on the linear $y$-axis versus exchange turn
    on the linear $x$-axis ($0$ to $30$) (steel-blue, positive region shaded
    amber), each forced act marked (crimson) and the closing line $D=0$
    marked (sea-green, dashed). A further act is committed exactly while
    $D>0$ and the exchange closes only at $D=0$, the empirical content of
    \Cref{thm:forcedact}: inquiry is forced by the gap.
    (\textbf{D})~Three-dimensional drift surface $1-\prod_i g_i$ over number
    of relays $n$ ($x$-axis, $1$ to $14$) and per-relay gain $g$ ($y$-axis,
    $0.4$ to $0.98$) (viridis colourmap), viewing angle elevation
    $26^{\circ}$ azimuth $-60^{\circ}$. Drift grows with chain length and
    with per-relay reshaping, the geometric picture of catalytic drift.}
    \label{fig:mechanics-4}
\end{figure*}

% =====================================================================
\section{Inquiry: a further act forced by the residual gap}
\label{sec:inquiry}
% =====================================================================

Finally we show that, because no report resolves to a held value, an
exchange cannot close on a value, so a further act is required; and that
this act is selected by the residual gap, not by any represented intent to
inquire. What an observer names a question is this forced act.

\begin{definition}[Exchange; residual gap]
\label{def:exchange}
An \emph{exchange} between two agents is an alternating sequence of accounts,
each agent's next account produced by a search seeded by the other's last.
The \emph{residual gap} at a stage is the standing cross-demand
$\dem>0$---the individuating contrast one side's account carries that the
other's search has not yet matched.
\end{definition}

\begin{theorem}[A further act is forced by the gap]
\label{thm:forcedact}
In an exchange, because no account resolves the queried item to a held value
(\Cref{ax:noncompletable,rem:noncompletable-role,thm:receiverrel}), the
exchange cannot terminate at a
value; while the residual gap $\dem>0$ persists, a further act is required to
reduce it, and that act is a further account. Were a value held, no further
act would be required; its unavailability forces the next act.
\end{theorem}

\begin{proof}
Termination at a value would require some account to resolve the item to a
residue-free point; none does (\Cref{ax:noncompletable,rem:noncompletable-role}),
and no shared value
is registered (\Cref{thm:receiverrel}). So the exchange cannot close on a
value. A positive residual gap $\dem>0$ is, by \Cref{def:exchange}, an
unmatched individuating contrast; reducing it requires an act that carries
the matching individuation---a further account (\Cref{thm:recsearch}). Were
a value held, the gap would be zero and no act required; hence the next act
is forced exactly by $\dem>0$.
\end{proof}

\begin{theorem}[The forced act is selected by the gap, not by intent]
\label{thm:gapselects}
The content of the forced further act is determined by the residual gap---by
which individuating contrast is still unmatched---and not by any represented
intent to inquire. The agent selects its next account as the search that
most reduces $\dem$; nowhere does it represent ``I am asking a question.''
\end{theorem}

\begin{proof}
The next act is the search that reduces the standing $\dem$
(\Cref{ax:purpsearch,def:exchange}); its target is fixed by where $\dem$ is
positive---the unmatched contrast---so the act's content is a function of the
gap. No term in this selection refers to a represented category ``question'':
the agent runs a $\dem$-descending search, the same operation as any report
(\Cref{thm:recsearch}). Hence the act is gap-selected and intent-free.
\end{proof}

So far ``forced by the gap'' applies to \emph{every} $\dem$-reducing act,
which is every act (\Cref{thm:recsearch,ax:purpsearch}); on its own it does
not single out inquiry. To make the result about \emph{questions}
non-vacuous we must say what distinguishes a question-act from a report-act,
and it is not intent---it is \emph{direction} on the residual gap.

\begin{definition}[Report-act and question-act, by direction on the gap]
\label{def:qdirection}
In an exchange, classify each act by its effect on the residual gap between
the two sides. A \emph{report-act} \emph{reduces} a gap the acting agent was
handed: it answers an incoming account, carrying the individuation that
matches an unmatched contrast the interlocutor's last account raised---it
lowers $\dem$ on the acting side. A \emph{question-act} \emph{raises} a gap
on the other side: it emits an account carrying an unmatched contrast the
interlocutor's next search must reduce---it creates $\dem$ in the
interlocutor. The two are opposite orientations of the same
$\dem$-descending search: one closes a received gap, the other opens one
across the exchange.
\end{definition}

\begin{theorem}[The forced act an observer names a question is the gap-opening one]
\label{thm:questionis}
The gap-forced further act of \Cref{thm:forcedact} is classified by an
observer as a \emph{question} exactly when it is a question-act in the sense
of \Cref{def:qdirection}: when it emits an unmatched contrast that opens a
residual gap in the interlocutor rather than closing one the agent received.
This classification is a property of the act's direction on the shared gap,
computed across the two agents, and refers to no represented intent to
inquire in the producing agent.
\end{theorem}

\begin{proof}
By \Cref{thm:forcedact} a positive gap forces a further $\dem$-descending
act, and by \Cref{thm:gapselects} its content is fixed by the gap, not by
intent. Whether an observer calls that act a report or a question turns on
its direction (\Cref{def:qdirection}): if it matches a contrast the agent was
handed, lowering the agent's own $\dem$, the observer reads a report; if it
carries a fresh unmatched contrast that the interlocutor's next search must
reduce, raising $\dem$ on the other side, the observer reads a question. This
is a relation between the two agents' successive gaps, decidable from the
accounts exchanged, with no term referring to a represented category in
either agent. Hence the specific result about inquiry: a question is the
gap-\emph{opening} forced act, named by the observer from its direction, not
a $\dem$-reducing act in general.
\end{proof}

\begin{corollary}[A question is a forced act, named by the observer]
\label{cor:question}
The act an observer names a question is the gap-forced further account of
\Cref{thm:forcedact}, selected by the residual gap (\Cref{thm:gapselects}),
and specifically the gap-\emph{opening} act of \Cref{thm:questionis}. The
agent represents no intent to inquire; ``question'' names the observer's
reading of the direction of a $\dem$-acting search, not a category the agent
holds.
\end{corollary}

\begin{proof}
By \Cref{thm:forcedact} the act is forced by $\dem>0$ and by
\Cref{thm:gapselects} its content is gap-selected with no represented intent;
the name ``question'' is the observer's classification of such an act.
\end{proof}

\begin{remark}[Asking without a concept of asking]
\label{rem:asking}
An agent of this kind asks---emits acts an interlocutor answers as
questions---without any concept of questioning, for the same reason it lies
without a concept of lying: the act is its ordinary purpose-directed search,
forced here by an unmatched contrast rather than initiated by an intent. A
population of such agents sustains inquiry with no agent that intends to
inquire; the inquiry is real on the observer's side and definition-free on
the agent's. We use only the theorems.
\end{remark}

% =====================================================================
\section{Where each name lives}
\label{sec:vocabulary}
% =====================================================================

We collect the paper's finding about the classifying vocabulary. In each
case the proof, not a premise, has placed the name. First we must define the
one term used informally up to now.

\begin{definition}[Testimony]
\label{def:testimony}
\emph{Testimony} is the act of emitting a report---a search terminus
(\Cref{def:searchprocess})---in response to a query. It is the only named
act in the observer's vocabulary that attaches to an agent-output, because
it is the only thing the agent does that is observable from outside: the
search interior is path-opaque (\Cref{thm:opacity,cor:noreconstruct}), so
nothing of the walk that produced the report is available to an observer,
only the report itself. Testimony is thus the emitted account, and nothing
more---it names the output, not any intent behind it.
\end{definition}

\begin{theorem}[Locus of the classifying vocabulary]
\label{thm:locus}
For an embedded agent under \Cref{ax:distfloor,ax:purpsearch,ax:actfloor,%
ax:noncompletable}:
\begin{enumerate}[label=\textup{(\roman*)},leftmargin=2.2em]
\item \textbf{Testimony} names an account $\acct$---a search terminus
emitted in response to a query (\Cref{def:testimony}, \Cref{thm:recsearch});
it is the one item of vocabulary that does attach to an agent-output, but it
attaches to the \emph{output}, not to any intent, and carries no commitment
to truth. It is non-trivially the only such term precisely because the
output is all that is observable, the interior being path-opaque
(\Cref{cor:noreconstruct}).
\item \textbf{Lie / false report} names the relation ``individuates a
different resting cut from the standard'' on a pair $(\acct,\acct^\ast)$
(\Cref{thm:falsityrelation}); it is undefined inside the producing agent
(\Cref{thm:falsityundefined}) and lives on the classifier.
\item \textbf{Rumour} names the multiplicative drift $1-G_n$ of a relayed
account (\Cref{thm:drift,cor:rumour}); it is a property of a chain of
relays, computed across agents, authored by none.
\item \textbf{Question} names a gap-forced, gap-\emph{opening} act
(\Cref{thm:forcedact,thm:gapselects,thm:questionis}); it lives on the
observer's reading of the act's direction on the shared gap, not on a
represented intent.
\end{enumerate}
In no case does the name denote a state or act the producing agent
represents as such.
\end{theorem}

\begin{proof}
Each clause restates a proved result: (i) \Cref{def:testimony,thm:recsearch};
(ii) \Cref{thm:falsityundefined,thm:falsityrelation}; (iii)
\Cref{thm:drift,cor:rumour}; (iv)
\Cref{thm:forcedact,thm:gapselects,thm:questionis}. In each, the classifying
property is shown to be either an output (testimony) or a relation computed
off outputs (lie, rumour, question), never an agent-internal representation
of the named category.
\end{proof}

\begin{corollary}[The agent classifies none of its own acts by these names]
\label{cor:agentnovocab}
The agent applies none of the terms \emph{testimony}, \emph{lie},
\emph{rumour}, or \emph{question} to its own acts. Each classification
requires data the agent's state does not hold: lying requires a standard
$\acct^\ast$ (absent, \Cref{thm:falsityundefined}, T4); rumour requires
tracking a chain of relays across agents (not in any one agent's state,
\Cref{thm:drift}, T5); a question requires representing an intent to inquire
or the interlocutor's gap direction (absent, \Cref{thm:gapselects,thm:questionis},
T6); and even testimony, as a \emph{named} category, requires the
outside standpoint from which the output is the only visible thing
(\Cref{def:testimony}). The agent's self-description, if it has one, is
exhausted by its current disposition and committed count
(\Cref{def:agentstate}); every one of these names is the observer's addition.
\end{corollary}

\begin{proof}
By T4 the agent holds no standard, so cannot apply ``lie''
(\Cref{thm:falsityundefined}); by T5 no single agent holds the relay chain,
so cannot apply ``rumour'' (\Cref{thm:drift}); by T6 the agent represents no
intent and computes no cross-agent gap direction, so cannot apply
``question'' (\Cref{thm:gapselects,thm:questionis}); and the agent's state is
the terminus-plus-count of \Cref{def:agentstate}, which contains none of
these predicates. Hence the agent classifies none of its own acts under the
observer's vocabulary.
\end{proof}

\begin{principle}[The vocabulary is the observer's]
\label{prin:observer}
The words with which we describe what embedded agents do to one another---
they testified, lied, gossiped, questioned---are, with the single partial
exception of ``testimony'' (which names an output), predicates the observer
computes on relations between accounts. The agents carry none of them. A
faithful model of social behaviour among such agents therefore need
implement no module for lying, gossip, or inquiry: these are what
purpose-directed, cost-bounded search in a finite world \emph{looks like}
from outside, and they appear whether or not anyone builds them in.
\end{principle}

% =====================================================================
\section{Numerical study}
\label{sec:validation}
% =====================================================================

The results are proved from the axioms; we exhibit them on explicit agents
to display their shape and check the derivations, computing every minimum cut
exactly by maximum flow.

\subsection{Method}

An agent is realised as a finite connected weighted world-graph with a
medium vertex adjacent to all items and all weights $\ge\flo$, a strongly
convex drive with known purpose, and an act-cost floor. On each agent we
compute, by exact enumeration and max-flow: separation costs and resting
cuts; search processes as $\dem$-descending walks and their reports; the
interiors of same-endpoint searches; registrations of a common event across
agents with different graphs; the falsity-relation on pairs of accounts under
a chosen standard; the search gains and the drift product $G_n$ along relay
chains; and the residual gap along an exchange with the forced-act selection.

\subsection{Predictions checked}

\begin{enumerate}[label=\textbf{(V\arabic*)},leftmargin=*]
\item \textbf{Recalling is searching.} Every report is a search terminus;
no report is returned without a committed act (\Cref{thm:recsearch}).
\item \textbf{Path opacity.} Same-seed, same-report searches have differing
interiors, and the report does not determine the interior
(\Cref{thm:opacity}); the committed record is strictly monotone
(\Cref{cor:noreconstruct}).
\item \textbf{Divergence without error.} A common event registers different
resting cuts across agents, each a correct min-cut individuation in its own
graph (\Cref{thm:receiverrel,cor:divergence}).
\item \textbf{Falsity is a relation.} The falsity-predicate has no value from
a single agent's state and a decidable value on a pair with a chosen standard
(\Cref{thm:falsityundefined,thm:falsityrelation}); a ``false'' and a ``true''
account are indistinguishable in producing-agent state
(\Cref{cor:noliar}).
\item \textbf{Multiplicative drift.} Surviving fraction along a relay chain
equals $\prod_i g_i$, decreasing geometrically for $g_i\le g<1$
(\Cref{thm:drift}).
\item \textbf{Forced inquiry.} An exchange never closes on a value; a further
act is committed exactly while $\dem>0$, its target the unmatched contrast
(\Cref{thm:forcedact,thm:gapselects}).
\end{enumerate}

\subsection{Results}

The suite is deterministic given its seed ($20260719$), runs a floor
$\flo=1$, and computes every minimum cut exactly by Edmonds--Karp
max-flow---no approximation enters. Across the nine prediction groups it
performs $740$ individual checks and \emph{all $740$ pass};
\Cref{tab:embedded-results} records the group-level outcomes.

\begin{table}[t]
\centering
\small
\caption{Numerical validation of \emph{Propagation Mechanics of an Embedded
Agent}. Every one of $740$ individual checks passes (seed $20260719$,
floor $\flo=1$, exact max-flow min-cut backend).}
\label{tab:embedded-results}
\begin{tabular}{@{}llr@{}}
\toprule
Prediction (theorem) & Checks & Margin / outcome \\
\midrule
V1 Floor forced (\Cref{thm:floorforced})   & $259$ & $\min\str/\flo = 1.092$ \\
V2 Recall is search (\Cref{thm:recsearch}) & $80$  & $\ge1$ act, no fetch \\
V3 Path opacity (\Cref{thm:opacity})       & $60$  & $60$ multi-interior \\
V4 Monotone record (\Cref{cor:noreconstruct}) & $60$ & record strictly $\uparrow$ \\
V5 Receiver-relative (\Cref{thm:receiverrel}) & $60$ & $60$ divergent, each correct \\
V6 Falsity is a relation (\Cref{thm:falsityrelation}) & $60$ & undefined in producer \\
V7 Catalytic drift (\Cref{thm:drift})      & $80$  & err.\ $=0$ \\
V8 Forced inquiry (\Cref{thm:forcedact})   & $80$  & targets argmax gap \\
V9 Locus of vocabulary (\Cref{thm:locus})  & $1$   & names placed \\
\midrule
\textbf{Total} & $\mathbf{740}$ & \textbf{all pass} \\
\bottomrule
\end{tabular}
\end{table}

The margins locate the mechanism. In V1 the smallest separation cost found
anywhere is $1.092\,\flo$: no item is ever individuated for free, and the
floor of \Cref{thm:floorforced} is close to tight, so the residual the whole
argument turns on is real and small, not a large safety margin. V3 and V5 are
the structural core: in all $60$ path-opacity graphs a report reached by
searches with distinct interiors admits more than one route, so the report
does not determine the interior; and in all $60$ divergence graphs a single
event registers a \emph{different} resting cut in each agent while each cut is
the exact min-cut individuation in its own graph---divergence with no error in
either. V6 confirms the paper's thesis operationally: the falsity-predicate
returns no value from any single agent's state and a decidable value only on a
pair of accounts with an elected standard, so a ``true'' and a ``false''
account are identical in producing-agent state. V7's relay drift matches
$\prod_i g_i$ to zero error, and V8's forced act targets the argmax residual
gap in every exchange.
\Cref{fig:mechanics-1,fig:mechanics-2,fig:mechanics-3,fig:mechanics-4}
display the floor, opacity, divergence, drift, and forced inquiry on the
explicit agents.




% =====================================================================
\section{Discussion}
\label{sec:discussion}
% =====================================================================

\subsection{What has and has not been established}

We have shown that an embedded agent with a purpose, a positive cost on its
acts, and access to its world only by search---set in a world it cannot
finish individuating---produces outputs an observer classifies as testimony,
false report, rumour, and questioning, and that it does so \emph{without
representing any of these}. We located each classifying name and found, as a
result rather than a premise, that---except for ``testimony,'' which names an
output---each is a relation an observer or second party computes between
accounts, undefined inside the producing agent. The agent that ``lies''
represents no lie; the relays that ``spread a rumour'' each merely searched;
the agent that ``asks'' holds no intent to inquire.

We have \emph{not} proved that such an agent must, in a given world, produce
any particular lie, rumour, or question: that is a claim of necessity, which
would require assumptions about the world's supply of engagements and
disagreements that we did not make. Our claim is emergence---these outcomes
arise without being represented---and it is the claim the axioms support.
Nor have we modelled how an observer chooses a standard $\acct^\ast$; the
falsity-relation's value depends on that choice, which is the classifier's,
not the agent's.

\subsection{Consequences for building embedded agents}

\Cref{prin:observer} has a direct design reading: a synthetic world of such
agents needs no subsystem for deception, gossip, or inquiry. Give each agent
a purpose, a cost on its acts, and search-only access to its own world-graph,
and let world-graphs differ across agents; then divergent testimony, drifting
rumour, and forced questioning appear as what that machinery looks like from
outside. Attempting to implement ``lying'' as a represented act would, by
\Cref{thm:falsityundefined}, be implementing something the agent has no place
to hold---and would be redundant with what search already produces.

\subsection{The social completion of the NPC architecture}
\label{sec:npc}

The results of this paper are the social-epistemological consequences of the
non-player-character architecture developed in the companion accounts. That
architecture equips an embedded agent with four implementation
invariants: a \emph{conserved identity} (a character invariant fixed under
re-description), a \emph{monotone committed count} (a history that only
grows), \emph{search-not-fetch} access (every report freshly searched,
\Cref{thm:recsearch,thm:fetchincoherent}), and \emph{phase exclusion}
(construction and commitment never coincide). An agent implemented under
these four invariants \emph{automatically} exhibits the behaviours this
paper characterises. Lying without a liar (\Cref{cor:noliar}), rumour
without an author (\Cref{cor:rumour}), questions without inquiry
(\Cref{cor:question})---these are not bugs to be prevented nor features to be
designed. They are what the four invariants \emph{look like from outside}
when several such agents interact. The observer's vocabulary---lie, rumour,
question---is the social-level description of what the individual-level
architecture produces; the monotone count makes fetch incoherent
(\Cref{thm:fetchincoherent}) and past states unreachable, search-not-fetch
makes every account route-opaque (\Cref{cor:noreconstruct}), and
receiver-relativity across agents makes divergence and drift inevitable
(\Cref{thm:receiverrel,thm:drift}). This paper is therefore not a standalone
result about embedded agents but the \emph{social completion} of the NPC
series: given the individual architecture, the social phenomena follow as
theorems.

\subsection{Relation to the companion accounts}

Two companion developments supply what this paper assumes. One derives why an
agent has an identity and a purpose at all---the conserved identity and the
attracting purpose; the other derives when an incoming interaction is
relevant and what a response costs---the bounded attention, the response
floor, and phase exclusion. The present paper takes purpose and act-cost as
given and asks what the agent \emph{does} when it answers the world---and
finds that the doing, run over a finite world by search under those same
invariants, is already testimony, divergence, drift, and inquiry, with the
names supplied entirely by whoever watches.

\subsection{Closing}

An embedded agent does not lie by deciding to, gossip by intending to, or ask
by meaning to. It searches its finite world toward its purpose at a cost, and
what comes out is classified---correctly, and only from outside---as lie,
rumour, or question. The behaviours do not wait on the concepts; the concepts
are names an observer brings to outputs the agent produces without them. None
of this needs to be sold: it is a set of theorems about search on finite
weighted graphs, and the readings of them as testimony and inquiry are
offered only as orientation.

\vspace{0.6em}
\hrule
\vspace{0.4em}
\noindent\emph{The agent searches; the observer names. Lie, rumour, and
question are relations between accounts, not acts of the agent that gave
them.}

\bibliographystyle{plain}
\bibliography{references}

\end{document}
